\section{Reproducibility and Ethical Considerations}\label{reproducibility-and-ethical-considerations}

The implementation uses PostgreSQL for the imported TAWOS source and Python modules for reconstruction, fitting, evaluation, and analysis. FastAPI and Docker Compose provide the surrounding local data-access stack; no production Agentick integration or practitioner pilot was conducted. The experimental runtime is Python 3.12.10, CatBoost 1.2.10, scikit-learn 1.9.1, pandas 3.0.6, and NumPy 2.5.3, with complete installed-version manifests retained.

The local record includes the pre-fit protocol/configuration/splits, inclusion and exclusion tables, feature dictionary, sampled SQL audits, model and calibration artifacts, test predictions, and per-sprint/project result tables. Six methods over four landmarks, fourteen folds, and two evaluation settings yield 672 method--landmark--fold artifact triplets; this includes deterministic baselines and should not be interpreted as 672 independently trained learning algorithms. Every triplet was checked against the frozen test cohort and sampled re-inference. Fourteen automated tests cover key replay, split, metric, and budget invariants. Tests support implementation correctness, not the external business validity of labels.

The underlying dataset is provided by its original authors \cite{tawosi2022tawos,tawosv11}. Its research-use terms and Apache 2.0 license are recorded in the data card. We do not use contributor identity as a feature or attempt re-identification. The artifact package currently resides in this local project and Git history; \textbf{no public repository, permanent artifact DOI, or independent reproduction is claimed}. The exact commands and paths are provided in Appendix A. Publication would require an appropriately prepared public package consistent with source terms.
